\subsection{P-12: an explicit width-three algebraic point and trace-coordinate span map}
\paragraph{Statement P-12.}
Let $n,D\ge1$.  Let $\mathcal C_{3,n,D}$ consist of the rational-coefficient polynomials of total degree at most $D$ having a width-at-most-three ROABP in some variable order.  The order is not given.  Put
\[
 B=nD+1,\qquad N=B^{21},\qquad b=(n-1)!,\qquad
 a=(n-1)B^{20}+B^{19}.
\]
For the nodes $0,\ldots,n-1$, let $\lambda_i$ be the Lagrange polynomial for node $i-1$, and define the explicit integer polynomials
\[
 A_i(T)=bT+
    \sum_{j=1}^{10}T^{B^{2j-1}}\,
       b\lambda_i\bigl(T^{B^{2j}}\bigr),
       \qquad 1\le i\le n.
 \tag{P12.1}
\]
Let $\theta$ be any root of $T^N-2$ and set $\alpha_i=A_i(\theta)$.  Then
\[
 0\ne f\in\mathcal C_{3,n,D}\quad\Longrightarrow\quad
 f(\alpha_1,\ldots,\alpha_n)\ne0.
 \tag{P12.2}
\]
The common number field $E=\mathbb Q(\theta)$ has degree $N$, and
\[
 \log|\Delta_E|\le N\log N+(N-1)\log2.
 \tag{P12.3}
\]
Every $A_i$ has degree at most $a$, at most $10n+1$ nonzero terms, and integer coefficients of bit length $O(n\log(n+1))$.  These data are constructible in time polynomial in $n,D$; the polynomial exponents are deliberately not optimized.

There is also an explicit one-parameter zero bound, valid over every characteristic-zero coefficient field:
\[
 f\ne0\quad\Longrightarrow\quad
 0\ne f(A_1(T),\ldots,A_n(T)),\qquad
 \deg f(A(T))\le Da<N.
 \tag{P12.4}
\]
Consequently the integer points $A(t)$ for $0\le t\le Da$ form a hitting set of size at most $Da+1$, with polynomial coordinate bit lengths.  This is a concrete construction for width three, including arbitrarily many rank-two cuts.

For a constant-boundary row product
\[
 R(x)=uA^{\rm mat}_1(x_{\sigma(1)})\cdots
          A^{\rm mat}_n(x_{\sigma(n)})\in\mathbb Q[x]^{1\times3}
\]
whose entries have total degree at most $D$, let $V\subseteq\mathbb Q^3$ be its coefficient-row span.  Write
\[
 R(\alpha)=\sum_{k=0}^{N-1}\theta^k r_k.
\]
Then the explicit arithmetic evaluation preserves that span in the following exact sense:
\[
 V=\operatorname{span}_{\mathbb Q}\{r_k:0\le k<N\}
  =\operatorname{span}_{\mathbb Q}
       \{\operatorname{Tr}_{E/\mathbb Q}(\beta_k R(\alpha)):0\le k<N\},
 \tag{P12.5}
\]
where
\[
 \beta_0=1/N,\qquad \beta_k=\theta^{N-k}/(2N)\quad(1\le k<N).
\]
The trace is applied entrywise.  The assertion concerns rational coefficient spans after expansion in the field basis, not the span of a single row over $E$.

\paragraph{Proof.}
The complete width-three rank-preservation proof in I-21 supplies the generator
\[
 G_i(t,\mathbf z,\mathbf y)=t+\sum_{j=1}^{10}z_j\lambda_i(y_j)
\]
for every variable order.  This is the specialized construction independently derived and checked from Liu's unrefereed rank-profile mechanism (2026, \S3.1, Theorem~3.2/Lemma~3.3 and \S3.2, Lemmas~3.4/3.6; P-B18).  I-21, rather than an assumption about the entire preprint, is the mathematical input here.

For a nonzero $f$, first put $f_b(x)=f(bx_1,\ldots,bx_n)$.  The nonzero scalar $b$ preserves nonvanishing, degree, and the width-three representation.  Hence $f_b(G)$ is a nonzero polynomial in its twenty-one parameters.  Its degree in $t$ and in each amplitude $z_j$ is at most $D$, and its degree in each selector $y_j$ is at most $(n-1)D$.  All these individual degrees are strictly below $B$.

Order the parameters as $t,z_1,y_1,\ldots,z_{10},y_{10}$ and substitute $T^{B^k}$ for the parameter in position $k$, starting with position zero.  Distinct parameter monomials have distinct resulting exponents, by uniqueness of base-$B$ expansion.  The substitution is therefore injective on the bounded-degree coefficient space.  Its value on $f_b(G)$ is precisely $f(A_1(T),\ldots,A_n(T))$, proving nonvanishing in (P12.4).  The same argument works over any characteristic-zero coefficient field.

Formula (P12.1) gives degree at most $a$ for every coordinate, so the composition has degree at most $Da$.  Explicitly,
\[
 Da=B^{20}\bigl((n-1)D+D/B\bigr)
      <B^{20}(nD)<B^{21}=N.
\]
The polynomial $T^N-2$ is Eisenstein at two and is consequently the minimal polynomial of $\theta$ over $\mathbb Q$.  A nonzero rational polynomial of degree below $N$ cannot vanish at $\theta$, proving (P12.2).

The Lagrange denominator at node $i-1$ is
$(-1)^{n-i}(i-1)!(n-i)!$, which divides $(n-1)!$.  Thus $b\lambda_i$ is integral.  Expansion of its linear factors bounds every coefficient by $(n!)^2$, a sufficient uniform bound including $n=1$; there are at most $n$ coefficients.  Distinct summands in (P12.1) have disjoint exponent patterns, except for zero coefficients, and give at most $10n+1$ terms.  This proves the stated degree, sparsity, and coefficient-size estimates.  Evaluation at the integers through $Da$ has coordinate bit length
\[
 O\bigl(n\log(n+1)+\log(10n+1)+a\log(Da+2)\bigr),
\]
which is polynomial in $n,D$.  The integer hitting claim follows from the univariate zero bound.

The order $\mathbb Z[\theta]$ has discriminant of absolute value
$N^N2^{N-1}$, obtained by multiplying the derivative values
$N\theta_j^{N-1}$ over the $N$ roots of $T^N-2$.  The field discriminant divides the order discriminant by the square of the order index.  This proves (P12.3).  No discriminant estimate is needed for the nonvanishing proof itself.

For the row-span assertion, the coefficients $r_k$ belong to $V$, because substitution and reduction in a rational basis take rational linear combinations of the original coefficient rows.  In fact the degree bound below $N$ means that no reduction is needed.  If their span were a proper subspace of $V$, there would be a rational column $v$ annihilating every $r_k$ but not every coefficient row of $R$.  The scalar polynomial $Rv$ would then be a nonzero member of $\mathcal C_{3,n,D}$ with $(Rv)(\alpha)=0$, contradicting (P12.2).  This proves the first equality in (P12.5).

Finally $\operatorname{Tr}_{E/\mathbb Q}(\theta^j)=0$ when $j$ is not divisible by $N$, while the traces of $1$ and $\theta^N$ are $N$ and $2N$.  For $0\le j,k<N$ this gives
$\operatorname{Tr}(\beta_k\theta^j)=\delta_{kj}$ directly.  Thus the displayed $\beta_k$ are the trace-dual basis, and the traces in (P12.5) equal the coefficient rows $r_k$.  $\square$

The arithmetic trace here is the ordinary number-field trace.  This construction supplies the previously missing effective algebraic point for the width-three class through rank preservation.  It does not infer that point from Frobenius equidistribution or provide a bound polynomial in a variable width.  Evaluating a black box at $\alpha$ requires access over $E$; the integer curve points in (P12.4) give an alternative using rational or integer queries.  Quantitative positive coefficient-energy equivalence remains subject to P-09.
